% Transcrição da ficha e do aviso de direitos autorais do PDF original.
\begin{pretextual}{}
\vspace*{2cm}
\noindent
\begin{minipage}{\linewidth}
\small
Zago, Natanael Henrik

\hspace*{1em}Usando Esquema GraphQL Para Geração De Consultas De Forma Aleatória / Natanael Henrik Zago. -- 2023.

\hspace*{1em}33 f.: il.

\hspace*{1em}Orientador: Prof. Dr. Samuel da Silva Feitosa.

\hspace*{1em}Trabalho de conclusão de curso (graduação) -- Universidade Federal da Fronteira Sul, curso de Ciência da Computação, Chapecó, SC, 2023.

\hspace*{1em}1. Geração de Código Aleatório. 2. GraphQL. 3. Testes de API.

I. Feitosa, Prof. Dr. Samuel da Silva, orientador. II. Universidade Federal da Fronteira Sul. III. Título.
\end{minipage}
\vfill
\noindent\textcopyright\ 2023

\noindent Todos os direitos autorais reservados a Natanael Henrik Zago. A reprodução de partes ou do todo deste trabalho só poderá ser feita mediante a citação da fonte.

\noindent E-mail: \texttt{natanael.zago@estudante.uffs.edu.br}
\end{pretextual}
